
\maketitle

\begin{abstract}
When does an empirically adequate theory admit a value-definite
completion --- an assignment of simultaneous definite values to all
observables?  We show that the answer depends on the algebraic
structure of the observations.  For a directed system of Boolean
algebras carrying compatible charges, the Stone construction produces
an unconditional measure on the dual space, and every observable
admits a simultaneous definite value via ultrafilters.  The passage to
a probability measure on an underlying state space requires only
$\sigma$-additivity, and the choice of state space is unconstrained by
the algebra~\cite{mahon2026}.

For orthomodular lattices --- the algebraic framework for incompatible
observations --- the situation is structurally different.  The
McDonald--Bimb\'{o} dual space carries principal filters as
distinguished structure, and the Kochen--Specker theorem blocks
global value-definiteness.  Probabilistic realism (descent to a
probability measure) remains available via Gleason-type results, but
value-definite realism does not.

The dividing line is distributivity.  We identify a four-level
coherence hierarchy --- consistency, contextual coherence,
probabilistic coherence, value-definite coherence --- and show that
distributivity of the observation algebra is the exact condition under
which all four levels are mutually commensurable.
\end{abstract}

\begin{quote}\small
\textbf{2020 Mathematics Subject Classification.}
03G12 (primary); 06C15, 81P10, 60A10 (secondary).

\textbf{Keywords.}
orthomodular lattice, Stone duality, value-definiteness,
Kochen--Specker theorem, Gleason's theorem, empirical adequacy,
constructive empiricism, Boolean algebra, directed system,
$\sigma$-additivity.
\end{quote}

% ------------------------------------------------------------------
\section{Introduction}
\label{sec:intro}
% ------------------------------------------------------------------

A central question in the philosophy of science asks when an
empirically successful theory can be completed to a realist one ---
one that assigns definite properties to an underlying reality, not
merely adequate predictions to observable phenomena.  Van
Fraassen~\cite{vanFraassen1980} argues that science need only aim at
\emph{empirical adequacy}: agreement with all observable phenomena.
The scientific realist demands more --- an underlying state space
whose properties generate the observations.

This debate is usually conducted in philosophical terms.  We give it
a mathematical formulation and show that its resolution depends on
the algebraic structure of the observations.

\subsection*{Three kinds of realism}

We distinguish three nested positions, each a strengthening of the
previous:

\begin{enumerate}[label=(\roman*)]
  \item \textbf{Empirical adequacy (EA).}
    A measure on the dual space of the observation algebra that
    agrees with all observed charges.  No commitment to an
    underlying state space.

  \item \textbf{Probabilistic realism (PR).}
    A probability measure on an underlying state space
    $(\Omega, \sigma(\mathcal{C}))$ that generates the observations.
    Requires descent from the dual-space measure.

  \item \textbf{Value-definite realism (VDR).}
    Every observable simultaneously has a definite value.  A global
    dispersion-free state --- equivalently, a 2-valued homomorphism
    on the full observation algebra.
\end{enumerate}

\noindent These are nested: VDR implies PR implies EA.  The question
is which implications reverse, and what governs the answer.

\subsection*{The role of the observation algebra}

Consider two settings:

\smallskip
\noindent\emph{Compatible observations.}
When all observations can be performed simultaneously, the
observation algebra is a Boolean algebra --- a complemented
distributive lattice.  Every Boolean algebra admits 2-valued
homomorphisms (ultrafilters exist by Zorn's lemma).  In this
setting, all three positions are available and mutually
commensurable.  The passage from EA to PR requires
$\sigma$-additivity of the charges~\cite{mahon2026}; the further
passage to VDR is free.

\smallskip
\noindent\emph{Incompatible observations.}
When some observations cannot be performed simultaneously --- as in
quantum mechanics, contextual experimental design, or distributed
systems with interfering sensors --- the observation algebra is an
orthomodular lattice (OML), a complemented but non-distributive
lattice.  The Kochen--Specker theorem~\cite{KochenSpecker1967}
shows that for the prototypical case $L(H)$ with $\dim H \geq 3$,
no 2-valued homomorphism exists: VDR is blocked.  Probabilistic
realism remains available via Gleason's
theorem~\cite{Gleason1957}, and empirical adequacy is unconditional.

\smallskip
The dividing line is distributivity.

\subsection*{Contributions}

We develop the following:

\begin{enumerate}
  \item A unified framework based on directed systems of
    complemented lattices, with the Boolean and orthomodular cases
    as instances (\S\ref{sec:boolean}--\S\ref{sec:oml-dual}).

  \item A four-level coherence hierarchy --- consistency, contextual
    coherence, probabilistic coherence, value-definite coherence ---
    with the obstruction at each level identified
    (\S\ref{sec:hierarchy}).

  \item A structural theorem: distributivity of the observation
    algebra is the exact condition under which all four levels are
    mutually commensurable (\S\ref{sec:commensurability}).

  \item An analysis of the philosophical consequences: the
    realism/empiricism debate has different mathematical answers
    depending on the observation algebra (\S\ref{sec:discussion}).
\end{enumerate}

% ------------------------------------------------------------------
\section{The Boolean case}
\label{sec:boolean}
% ------------------------------------------------------------------

We recall the framework and results of~\cite{mahon2026}, which treats
directed systems of Boolean algebras.

\subsection*{Directed systems and compatible charges}

Let $(\iota, \leq)$ be a directed set indexing a family of Boolean
algebras $\{B_i\}_{i \in \iota}$ with injective homomorphisms
$\varphi_{ij} : B_i \to B_j$ for $i \leq j$.  Write
$\mathcal{C} = \bigcup_i B_i$ for the direct limit --- the
\emph{cylinder algebra}.

A \emph{compatible family of normalised charges} is a collection
$\{\ell_i\}$ with $\ell_i : B_i \to [0,1]$ finitely additive,
$\ell_i(1) = 1$, and $\ell_j \circ \varphi_{ij} = \ell_i$ for
$i \leq j$.

\subsection*{The Stone construction}

The Stone space $\St(\mathcal{C})$ is the space of ultrafilters on
$\mathcal{C}$, compact and totally disconnected.  Each
$\omega \in \Omega$ determines a principal ultrafilter
$\pure(\omega) = \{E \in \mathcal{C} : \omega \in E\}$.

Any compatible family of normalised charges extends to a
$\sigma$-additive Baire probability measure $\hat{\mu}$ on
$\St(\mathcal{C})$ --- unconditionally, with no assumption of
$\sigma$-additivity~\cite{choksi1958, mahon2026}.  This is the
\emph{empirically adequate} object: it agrees with all charges on
cylinder events.

\subsection*{Descent and $\sigma$-additivity}

The measure $\hat{\mu}$ lives on $\St(\mathcal{C})$, not on
$\Omega$.  Descent to a probability measure $P$ on
$(\Omega, \sigma(\mathcal{C}))$ requires
$\hat{\mu}(\pure(\Omega)) = 1$ --- that $\hat{\mu}$ concentrates on
the principal ultrafilters.  This holds if and only if each charge
$\ell_i$ is $\sigma$-additive~\cite{mahon2026}.

\subsection*{Value-definiteness in the Boolean case}

Every Boolean algebra admits ultrafilters (by Zorn's lemma), and
every ultrafilter is a 2-valued homomorphism.  A 2-valued
homomorphism assigns to each element $a \in \mathcal{C}$ a definite
value in $\{0, 1\}$, respecting the lattice operations.

Moreover, any subset $S \subseteq \St(\mathcal{C})$ that projects
surjectively onto each $\St(B_i)$ can serve as a realisation space:
set $\Omega = S$ with evaluation maps given by the restriction maps.
The algebra imposes no constraint on the choice of
$\Omega$~\cite{mahon2026}.

\begin{proposition}[Boolean commensurability]
\label{prop:boolean-comm}
For a directed system of Boolean algebras with compatible charges:
\begin{enumerate}[label=(\alph*)]
  \item EA is unconditional: $\hat{\mu}$ on $\St(\mathcal{C})$
    always exists.
  \item PR holds iff each $\ell_i$ is $\sigma$-additive.
  \item VDR holds unconditionally: $\St(\mathcal{C})$ consists
    entirely of 2-valued homomorphisms.
\end{enumerate}
In particular, EA, PR, and VDR are mutually commensurable: the
only obstruction is $\sigma$-additivity at the EA--PR transition,
and it is a condition on the charges, not the algebra.
\end{proposition}

\begin{proof}
Part~(a) is the Stone construction~\cite{choksi1958, mahon2026}.
Part~(b) is~\cite[Theorem~1]{mahon2026}.  Part~(c): ultrafilters on
a Boolean algebra are exactly the 2-valued homomorphisms; their
existence follows from Zorn's lemma applied to the filter generated
by any non-zero element.
\end{proof}

% ------------------------------------------------------------------
\section{Orthomodular observation algebras}
\label{sec:oml}
% ------------------------------------------------------------------

The Boolean assumption is that all observations are compatible: any
finite collection can be performed simultaneously.  When this fails
--- when some observations preclude others --- the observation
algebra loses distributivity.  The appropriate generalisation is the
orthomodular lattice.

\subsection*{Orthomodular lattices}

\begin{definition}
An \emph{ortholattice} is a bounded lattice
$A = (A,\, \wedge,\, \vee,\, {}^\perp,\, 0,\, 1)$ with an
orthocomplementation ${}^\perp : A \to A$ satisfying
$a \wedge a^\perp = 0$, \; $a \vee a^\perp = 1$, \;
$a^{\perp\perp} = a$, and \; $a \leq b \Rightarrow b^\perp \leq a^\perp$.
\end{definition}

\begin{definition}
An \emph{orthomodular lattice} (OML) is an ortholattice satisfying
the \emph{orthomodular law}:
\[
  a \leq b
  \quad\Longrightarrow\quad
  b = a \vee (a^\perp \wedge b).
\]
\end{definition}

\noindent Every Boolean algebra is an OML (the orthomodular law
follows from distributivity), but the converse fails.  An OML is
Boolean if and only if it is
distributive~\cite[Proposition~2.3]{McDonaldBimbo2023}.

\begin{example}
Let $H$ be a separable Hilbert space with $\dim H \geq 3$.  The
lattice $L(H)$ of closed linear subspaces, ordered by inclusion with
$U^\perp = \{x \in H : \langle x, y \rangle = 0 \;\forall y \in U\}$,
is an orthomodular lattice that is not
distributive~\cite{Kalmbach1983}.
\end{example}

\begin{example}[Incompatible experimental design]
Consider an experiment with three binary measurements $A$, $B$, $C$,
where $(A,B)$ and $(A,C)$ can be performed jointly but $(B,C)$
cannot.  The jointly measurable pairs generate Boolean subalgebras,
but the full observation algebra is orthomodular, not Boolean.
This situation arises in contextual behavioural science, clinical
trials with mutually exclusive protocols, and sensor networks with
interfering instruments.
\end{example}

\subsection*{Commutativity and classical contexts}

Two elements $a, b$ in an OML are said to \emph{commute} if they
generate a Boolean subalgebra.  A maximal set of pairwise commuting
elements forms a \emph{Boolean context} --- a maximal Boolean
subalgebra of $A$.  Within each context, classical logic and
classical probability apply.  The non-classical character of the
OML arises from the relationships \emph{between} contexts.

\subsection*{States on OMLs}

\begin{definition}
A \emph{state} on an OML $A$ is a function $s : A \to [0,1]$ with
$s(1) = 1$ and
\[
  a \perp b
  \quad\Longrightarrow\quad
  s(a \vee b) = s(a) + s(b),
\]
where $a \perp b$ means $a \leq b^\perp$.
\end{definition}

\noindent A state is additive on \emph{orthogonal} pairs only ---
not on arbitrary joins.  This is the OML analogue of a finitely
additive charge on a Boolean algebra.

\begin{definition}
A state $s$ is \emph{dispersion-free} if $s(a) \in \{0,1\}$ for
all $a \in A$.
\end{definition}

\noindent A dispersion-free state is a 2-valued homomorphism: it
assigns a definite truth value to every proposition in $A$.  On a
Boolean algebra, dispersion-free states are exactly the
ultrafilters.

\subsection*{The Kochen--Specker obstruction}

\begin{theorem}[Kochen--Specker~\cite{KochenSpecker1967}]
\label{thm:KS}
If $\dim H \geq 3$, there is no dispersion-free state on $L(H)$.
\end{theorem}

\noindent No assignment of simultaneous definite values to all
observables is consistent with the algebraic structure.  This is
not a limitation of knowledge but a structural impossibility: the
OML $L(H)$ does not admit a 2-valued homomorphism.

The obstruction is dimensional: for $\dim H = 2$, dispersion-free
states exist (every maximal Boolean subalgebra of $L(\mathbb{C}^2)$
is a four-element Boolean algebra, and 2-valued homomorphisms are
plentiful).  The critical threshold is $\dim H = 3$.

\subsection*{Gleason's theorem}

Despite the absence of dispersion-free states, \emph{probabilistic}
states on $L(H)$ are well-behaved:

\begin{theorem}[Gleason~\cite{Gleason1957}]
\label{thm:gleason}
If $\dim H \geq 3$, every $\sigma$-additive state on $L(H)$ is of
the form $s(P) = \mathrm{tr}(\rho P)$ for a unique density operator
$\rho$ on~$H$.
\end{theorem}

\noindent Gleason's theorem says that probabilistic realism is
available for $L(H)$: every state extends to a well-defined
probability assignment via the Born rule.  What fails is
value-definite realism --- the demand that every observable have a
simultaneous definite value.

% ------------------------------------------------------------------
\section{The OML dual space}
\label{sec:oml-dual}
% ------------------------------------------------------------------

In the Boolean case, Stone duality provides the dual space
$\St(\mathcal{C})$ of ultrafilters.  For OMLs, the analogue is the
McDonald--Bimb\'{o} duality~\cite{McDonaldBimbo2023}, which we now
describe.

\subsection*{Filters on an OML}

Let $A$ be an OML.  A \emph{filter} on $A$ is a non-empty subset
$x \subseteq A$ that is upward closed ($a \in x$ and $a \leq b$
implies $b \in x$) and closed under finite meets ($a, b \in x$
implies $a \wedge b \in x$).  A filter is \emph{proper} if
$0 \notin x$; the unique \emph{improper} filter is $\omega = A$.

For each $a \in A$, the \emph{principal filter} generated by $a$ is
${\uparrow}a = \{b \in A : a \leq b\}$.  Write $F(A)$ for the
collection of all filters on $A$, and $P(A) \subseteq F(A)$ for the
subcollection of principal filters.

\subsection*{The dual space}

\begin{definition}[McDonald--Bimb\'{o}~\cite{McDonaldBimbo2023}]
\label{def:oml-dual}
Let $A$ be an OML.  The \emph{dual space} of $A$ is the ordered
relational topological space
\[
  S_0(A) = \bigl(F(A),\; \subseteq,\; \perp_A,\; P(A),\; \mathcal{T}\bigr),
\]
where:
\begin{enumerate}[label=(\roman*)]
  \item $F(A)$ is the set of all filters on $A$;
  \item $\subseteq$ is set inclusion;
  \item $x \perp_A y$ iff there exists $a \in A$ with $a \in x$ and
    $a^\perp \in y$;
  \item $P(A)$ is the collection of principal filters;
  \item $\mathcal{T}$ is the topology generated by the subbasis
    $\{h(a) : a \in A\} \cup \{F(A) \setminus h(a) : a \in A\}$,
    where $h(a) = \{x \in F(A) : a \in x\}$.
\end{enumerate}
\end{definition}

\begin{theorem}[McDonald--Bimb\'{o}~\cite{McDonaldBimbo2023}]
\label{thm:oml-duality}
$S_0(A)$ is an orthomodular space (compact, totally disconnected).
The category of OMLs and homomorphisms is dually equivalent to the
category of orthomodular spaces and continuous weak p-morphisms.
Moreover, $A \cong \CO(S_0(A))^\dagger$, the OML of clopen
$\perp$-stable subsets of $S_0(A)$.
\end{theorem}

\subsection*{Structural comparison with Stone duality}

The critical differences are collected in the following table:

\medskip
\begin{center}
\begin{tabular}{@{}lll@{}}
\toprule
Feature & Boolean (Stone) & OML (McDonald--Bimb\'{o}) \\
\midrule
Dual points & Ultrafilters & All filters \\
Distinguished subset & None & $P(A)$ (principal filters) \\
Clopens & Boolean algebra & OML \\
Orthogonality & Trivial & Non-trivial ($\perp_A$) \\
2-valued homomorphisms & Plentiful (Zorn) & May not exist (KS) \\
\bottomrule
\end{tabular}
\end{center}
\medskip

\noindent In Stone duality, the principal filters (i.e., the
principal ultrafilters $\pure(\omega)$ for $\omega \in \Omega$) are
\emph{invisible}: they play no distinguished role in the dual space.
Which ultrafilters are principal depends on the choice of
$\Omega$~\cite{mahon2026}.

In the McDonald--Bimb\'{o} duality, $P(A)$ is \emph{part of the
structure} of the dual space: the orthomodular frame conditions
(Definition~\ref{def:oml-dual}) explicitly reference $P(A)$.  This
is the algebraic formalisation of the observation that, for
non-distributive algebras, the distinction between realised and
unrealised observation-patterns is not freely chosen but is
constrained by the algebra.

% ------------------------------------------------------------------
\section{Three kinds of realism, revisited}
\label{sec:three-kinds}
% ------------------------------------------------------------------

We now have the framework to state the three positions precisely in
both settings.

\begin{definition}
Let $A$ be a complemented lattice (Boolean or orthomodular) with a
state $s : A \to [0,1]$.
\begin{enumerate}[label=(\roman*)]
  \item \textbf{Empirical adequacy (EA):}
    $s$ defines a consistent assignment $\mu(h(a)) = s(a)$ on the
    clopens of the dual space.

  \item \textbf{Probabilistic realism (PR):}
    $\mu$ extends to a $\sigma$-additive probability measure on the
    Baire $\sigma$-algebra of the dual space, concentrating on
    $P(A)$.

  \item \textbf{Value-definite realism (VDR):}
    There exists a dispersion-free state on $A$ --- a 2-valued
    homomorphism assigning simultaneous definite values to all
    elements.
\end{enumerate}
\end{definition}

\begin{proposition}[OML status]
\label{prop:oml-status}
For $A = L(H)$ with $\dim H \geq 3$:
\begin{enumerate}[label=(\alph*)]
  \item EA holds: every state defines a consistent assignment on
    $\CO(S_0(A))^\dagger$.
  \item PR holds: Gleason's theorem~\textup{(\ref{thm:gleason})}
    gives a $\sigma$-additive measure via $s(P) = \mathrm{tr}(\rho P)$.
  \item VDR fails: the Kochen--Specker
    theorem~\textup{(\ref{thm:KS})} shows no dispersion-free state
    exists.
\end{enumerate}
\end{proposition}

\noindent Comparing with Proposition~\ref{prop:boolean-comm}:

\medskip
\begin{center}
\begin{tabular}{@{}lccc@{}}
\toprule
Position & Boolean & OML ($\dim \geq 3$) \\
\midrule
EA & \checkmark\, (unconditional) & \checkmark\, (unconditional) \\
PR & \checkmark\, (iff $\sigma$-additive) & \checkmark\, (Gleason) \\
VDR & \checkmark\, (ultrafilters) & $\times$\, (Kochen--Specker) \\
\bottomrule
\end{tabular}
\end{center}
\medskip

\noindent The single structural difference --- distributivity versus
non-distributivity --- controls exactly the VDR row.

% ------------------------------------------------------------------
\section{The descent problem for OMLs}
\label{sec:descent}
% ------------------------------------------------------------------

In the Boolean case, the descent from $\hat{\mu}$ on
$\St(\mathcal{C})$ to $P$ on $(\Omega, \sigma(\mathcal{C}))$
requires $\sigma$-additivity of the charges and a choice of
realisation space $\Omega$.  The algebra imposes no constraint on
$\Omega$.

For OMLs, the descent problem has a different character.  A state
$s$ on $A$ defines $\mu(h(a)) = s(a)$ on the clopen $\perp$-stable
sets $\CO(S_0(A))^\dagger$.

\subsection*{The extension obstacle}

In the Boolean case, $\Clop(\St(\mathcal{C}))$ is a Boolean algebra
of sets, and $s$ is finitely additive on it.  The
Carath\'{e}odory--Choksi extension machinery applies: compactness of
$\St(\mathcal{C})$ gives $\sigma$-subadditivity, and the extension
to a Baire measure is automatic.

In the OML case, $\CO(S_0(A))^\dagger$ is an OML of sets --- closed
under $\cap$ and ${}^\perp$ but \emph{not} under arbitrary $\cup$.
The state $s$ is additive on orthogonal pairs only.
Carath\'{e}odory does not apply: there is no Boolean algebra of
sets on which $s$ is finitely additive.

For the prototypical case $A = L(H)$ with $\dim H \geq 3$,
Gleason's theorem (Theorem~\ref{thm:gleason}) provides the
extension: every state arises from a density operator, giving a
$\sigma$-additive measure via the Born rule.  For general OMLs,
the extension problem is open.

\subsection*{Partial value-definite descent}

Even when probabilistic descent succeeds (Gleason), value-definite
descent does not.  The Bub--Clifton
theorem~\cite{BubClifton1996} addresses what \emph{can} be
salvaged:

\begin{theorem}[Bub--Clifton~\cite{BubClifton1996}]
\label{thm:bub-clifton}
Given a state $\rho$ on $L(H)$ and a preferred observable $R$,
there is a unique maximal sublattice
$D(\rho, R) \subseteq L(H)$ such that:
\begin{enumerate}[label=(\roman*)]
  \item $D(\rho, R)$ is a Boolean algebra;
  \item $D(\rho, R)$ contains the spectral projections of $R$;
  \item The restriction of $\rho$ to $D(\rho, R)$ is a mixture of
    dispersion-free states.
\end{enumerate}
\end{theorem}

\noindent In our terminology: $D(\rho, R)$ is the maximal Boolean
subalgebra on which VDR is available --- the largest ``classical
context'' in which all observables simultaneously have definite
values, given the state $\rho$ and the preferred observable $R$.

This is a \emph{partial} and \emph{context-dependent} descent.  It
is partial because $D(\rho, R) \subsetneq L(H)$: not all
observables receive definite values.  It is context-dependent
because $D(\rho, R)$ depends on the choice of $R$ --- different
preferred observables yield different maximal Boolean subalgebras.

The Bub--Clifton result thus sharpens the Kochen--Specker
obstruction: VDR fails globally, but it succeeds locally --- within
each Boolean context.  The failure is in the \emph{global}
assembly, not in any individual context.

% ------------------------------------------------------------------
\section{Distributivity as the commensurability condition}
\label{sec:commensurability}
% ------------------------------------------------------------------

We can now state the main structural result.

\begin{theorem}[Commensurability]
\label{thm:commensurability}
Let $A$ be a complemented lattice.
\begin{enumerate}[label=(\alph*)]
  \item If $A$ is distributive \textup{(}Boolean\textup{)}, then
    EA, PR, and VDR are mutually commensurable:
    \begin{itemize}
      \item EA is unconditional
        \textup{(}Stone construction\textup{)}.
      \item PR holds iff the charges are $\sigma$-additive
        \textup{\cite{mahon2026}}.
      \item VDR holds unconditionally
        \textup{(}ultrafilters exist\textup{)}.
    \end{itemize}
    Moreover, the choice of realisation space $\Omega$ is
    unconstrained by the algebra.

  \item If $A$ is non-distributive \textup{(}OML with
    $A \cong L(H)$, $\dim H \geq 3$\textup{)}, then:
    \begin{itemize}
      \item EA is unconditional.
      \item PR holds \textup{(}Gleason,
        Theorem~\textup{\ref{thm:gleason})}.
      \item VDR fails globally
        \textup{(}Kochen--Specker,
        Theorem~\textup{\ref{thm:KS})}.
      \item VDR holds locally within each Boolean context
        \textup{(}Bub--Clifton,
        Theorem~\textup{\ref{thm:bub-clifton})}.
    \end{itemize}
\end{enumerate}
\end{theorem}

\begin{proof}
Part~(a) is Proposition~\ref{prop:boolean-comm}.  For
part~(b): EA follows from the existence of states on $L(H)$ and the
McDonald--Bimb\'{o} dual~\cite{McDonaldBimbo2023}; PR from
Gleason's theorem; VDR-failure from Kochen--Specker;
local VDR from Bub--Clifton.
\end{proof}

\begin{corollary}
\label{cor:distributivity}
Distributivity of the observation algebra is the exact condition
under which value-definite realism and empirical adequacy are
commensurable.
\end{corollary}

\begin{proof}
If $A$ is distributive, VDR holds
(Theorem~\ref{thm:commensurability}(a)).  If $A$ is
non-distributive with $A \cong L(H)$ and $\dim H \geq 3$, VDR
fails (Theorem~\ref{thm:commensurability}(b)).  Since an OML is
Boolean iff it is distributive~\cite{McDonaldBimbo2023}, the
dividing line is exactly distributivity.
\end{proof}

\begin{remark}
The corollary requires the dimensional condition $\dim H \geq 3$
for the Kochen--Specker direction.  For $\dim H = 2$, the OML
$L(\mathbb{C}^2)$ is non-distributive but admits dispersion-free
states: VDR holds despite non-distributivity.  The sharp
characterisation is: VDR is commensurable with EA iff $A$ is
distributive \emph{or} every non-distributive component has
dimension at most $2$.  In physically relevant cases
($\dim H \geq 3$), distributivity is the effective dividing line.
\end{remark}

% ------------------------------------------------------------------
\section{The coherence hierarchy}
\label{sec:hierarchy}
% ------------------------------------------------------------------

The preceding analysis reveals a four-level hierarchy of coherence
conditions on an observation algebra $A$, each strictly stronger
than the last.

\begin{definition}[Coherence levels]
\label{def:hierarchy}
Let $A$ be a complemented lattice with a family of charges.
\begin{enumerate}
  \item[\textbf{Level~1.}] \textbf{Consistency.}
    No contradiction within any single Boolean context: each
    maximal Boolean subalgebra of $A$ admits a 2-valued
    homomorphism compatible with the charges.

  \item[\textbf{Level~2.}] \textbf{Contextual coherence.}
    The assignments across all Boolean contexts are compatible: a
    state $s : A \to [0,1]$ exists, agreeing with the charges on
    each context.

  \item[\textbf{Level~3.}] \textbf{Probabilistic coherence.}
    The state extends to a $\sigma$-additive probability measure on
    the dual space of $A$ --- i.e., PR holds.

  \item[\textbf{Level~4.}] \textbf{Value-definite coherence.}
    A global dispersion-free state exists: a 2-valued homomorphism
    on the full algebra $A$ --- i.e., VDR holds.
\end{enumerate}
\end{definition}

\noindent Each level implies the previous one:
$4 \Rightarrow 3 \Rightarrow 2 \Rightarrow 1$.  The converses
fail in general.  The following table records which transitions are
obstructed:

\medskip
\begin{center}
\begin{tabular}{@{}llll@{}}
\toprule
Transition & Boolean & OML ($\dim \geq 3$) & Obstruction \\
\midrule
$1 \to 2$ & Trivial & Non-trivial & Contextuality \\
$2 \to 3$ & $\sigma$-additivity & Gleason-type & Infinitary \\
$3 \to 4$ & Free (ultrafilters) & Blocked & Kochen--Specker \\
\bottomrule
\end{tabular}
\end{center}
\medskip

\noindent In the Boolean case, transitions $1 \to 2$ and $3 \to 4$
are trivial: distributivity collapses them.  The only non-trivial
transition is $2 \to 3$, which is the content
of~\cite{mahon2026}.

In the OML case, every transition is non-trivial.  The critical
obstruction is at $3 \to 4$: the Kochen--Specker theorem blocks the
passage from probabilistic coherence to value-definite coherence.
This is the precise level at which distributivity is load-bearing.

\begin{remark}[Relation to contextuality hierarchies]
Abramsky and Brandenburger~\cite{AbramskyBrandenburger2011}
introduced a hierarchy of contextuality: probabilistic,
possibilistic (logical), and strong contextuality.  Their hierarchy
grades the \emph{strength of the obstruction} to a global section of
a presheaf of probability distributions.  Our hierarchy grades
\emph{coherence demands}: what kind of global object the
observation algebra can support.  The two hierarchies are
complementary and orthogonal: theirs asks ``how badly does a global
section fail?''; ours asks ``what kind of global object can the
algebra support?''
\end{remark}

% ------------------------------------------------------------------
\section{Discussion}
\label{sec:discussion}
% ------------------------------------------------------------------

\subsection*{The realism debate is not uniform}

The standard philosophical debate between scientific realism and
constructive empiricism treats the question as uniform: either
theories aim at truth or they aim at empirical adequacy.  The
algebraic analysis above shows that the question is
\emph{parametrised} by the observation algebra.

For Boolean observations, the debate is genuine: the mathematics
does not prefer one position over another.  EA, PR, and VDR are all
available; $\sigma$-additivity is the only bridge condition, and it
is a property of the charges, not of the algebra.

For non-Boolean observations, the mathematics is not neutral.  VDR
is algebraically impossible (for $\dim \geq 3$); PR is available
but structurally constrained; EA is always available.  The
constructive empiricist position~\cite{vanFraassen1980} is the
generic one.  Value-definite realism is the special case --- the
case in which the observation algebra happens to be distributive.

\subsection*{Historical echoes}

The distinction among the three positions illuminates classical
debates:
\begin{itemize}
  \item \emph{Einstein's realism} (``God does not play dice'') is
    the demand for VDR.  The Kochen--Specker theorem shows this is
    not merely philosophically contestable but algebraically
    impossible for $\dim \geq 3$.

  \item \emph{Bohr's complementarity} is the observation that
    different Boolean contexts cannot be simultaneously realised.
    In our hierarchy, this is the non-triviality of the $1 \to 2$
    transition.

  \item \emph{Bell's beables}~\cite{HalvorsonClifton1999} ---
    the observables that \emph{can} have definite values --- are
    exactly the elements of the maximal Boolean subalgebra
    $D(\rho, R)$ identified by Bub--Clifton
    (Theorem~\ref{thm:bub-clifton}).

  \item \emph{De Finetti's operationalism} is EA in the Boolean
    case: probability as coherent betting rates on observables,
    without commitment to an underlying state space.
\end{itemize}

\subsection*{Beyond quantum mechanics}

The analysis applies wherever observations are incompatible ---
not only in quantum mechanics.  Orthomodular structure arises in
contextual behavioural science, distributed systems with
interfering sensors, database queries with access restrictions, and
experimental designs with mutually exclusive protocols.  In each
setting, the question ``does empirical adequacy entail
value-definiteness?'' has the same algebraic answer: only if the
observation algebra is distributive.

\subsection*{Open questions}

Several directions remain:
\begin{enumerate}
  \item \emph{The OML extension problem.}
    Does every state on a general OML extend to a $\sigma$-additive
    measure on $S_0(A)$, as in the Boolean case?  Gleason's theorem
    answers this for $L(H)$; the general case is open
    (\S\ref{sec:descent}).

  \item \emph{Intermediate algebras.}
    Between Boolean algebras and arbitrary OMLs lie the
    MO$_n$-free OMLs, modular ortholattices, and other
    subvarieties.  Where exactly does VDR fail in this spectrum?

  \item \emph{Dynamical structure.}
    If the observation algebra carries a group action (stationarity,
    exchangeability), how does symmetry interact with the coherence
    hierarchy?  The Boolean case is partially understood
    (de~Finetti's theorem for exchangeability); the OML case is
    unexplored.

  \item \emph{The topos connection.}
    The D\"{o}ring--Isham topos approach~\cite{DoeringIsham2008}
    addresses closely related questions via spectral presheaves
    rather than the McDonald--Bimb\'{o} filter-space duality.  A
    comparison of the two frameworks --- and whether the
    commensurability theorem has a topos-theoretic formulation ---
    is a natural next step.
\end{enumerate}

\ifstandalone
\bibliographystyle{plain}
\bibliography{references}
\fi
